\section*{ПРИЛОЖЕНИЕ Г. Практическая интеграция и инфраструктурный контекст}

\subsection*{Г.1. Место сервиса в общей продуктовой картине}

Разработанный микросервис проектировался не как абстрактный
исследовательский прототип, а как компонент продуктовой
инфраструктуры компании, в которой одновременно эксплуатируются
системы защиты от DDoS-атак (BIFIT MITIGATOR) и системы анализа
сетевого трафика (BIFIT COLLECTOR). Обе категории продуктов
полагаются на актуальное знание принадлежности IP-адресов автономным
системам, но используют его по-разному:

\begin{itemize}
\item для систем защиты ASN-метка является сигналом, на основе
которого принимаются решения о фильтрации, скоринге и распределении
нагрузки;
\item для систем анализа ASN-метка участвует в построении сводных
отчётов, в детектировании аномалий и в обогащении flow-записей.
\end{itemize}

Принципиально важно, что обоим продуктам нужен один и тот же ответ
на вопрос «к какой ASN принадлежит данный IP», но раньше каждый из
них извлекал этот ответ собственным способом, с собственным
жизненным циклом обновления и собственной моделью консистентности.

\subsection*{Г.2. Эффект унификации источника данных}

Появление общего сервиса даёт несколько системных эффектов, не
сводимых только к ускорению одной операции.

Во-первых, исчезает источник скрытого расхождения данных.
Когда два продукта строят свои таблицы независимо, их представления
о принадлежности одного и того же IP могут разойтись, что особенно
неприятно при разборе инцидентов, где требуется единая трактовка
наблюдений. Унификация источника данных снимает этот класс проблем
по построению.

Во-вторых, упрощается операционная модель. Поддерживать актуальной
одну подсистему дешевле, чем поддерживать актуальными несколько
независимых реализаций одной и той же функции. Это особенно заметно
при изменениях в источнике данных или в формате представления BGP.

В-третьих, появляется естественная точка для расширения. Например,
можно добавить в сервис способность возвращать дополнительные
маршрутизационные атрибуты, не меняя при этом архитектуру каждого
продукта-потребителя. С точки зрения evolvability это очень
существенный сдвиг по сравнению с разрозненной реализацией.

\subsection*{Г.3. Снижение совокупной стоимости владения}

Сравнительная экономия инфраструктурного контура от появления
общего сервиса заметна сразу по нескольким направлениям:

\begin{itemize}
\item на уровне CPU -- благодаря уменьшению числа повторных пересборок
таблицы в каждом продукте;
\item на уровне памяти -- благодаря отсутствию дублирующихся
in-memory представлений одного и того же состояния;
\item на уровне сети -- благодаря тому, что подключения к источнику
маршрутных данных консолидируются в одной точке;
\item на уровне эксплуатации -- благодаря единым средствам
наблюдаемости, профилирования и диагностики;
\item на уровне разработки -- благодаря единой кодовой базе для
поддержки маршрутизационного отображения.
\end{itemize}

Каждый из этих пунктов в отдельности можно рассматривать как чисто
инженерное улучшение, но в совокупности они формируют качественно
другую модель эксплуатации.

\subsection*{Г.4. Влияние на жизненный цикл продуктов}

Помимо непосредственной экономии, сервис влияет на скорость
эволюции продуктов. Раньше любое изменение в способе получения
ASN-данных приводило к одинаковым изменениям в нескольких кодовых
базах. Теперь оно локализовано в одном сервисе, а продукты-потребители
получают изменение «бесплатно» в рамках обычного обновления
зависимостей API.

Это особенно важно для следующих типичных сценариев:

\begin{itemize}
\item смена источника BGP (например, переход на другой узел GoBGP
или подключение дополнительного источника);
\item изменение модели обработки IPv6 (например, более тонкая
сегментация \texttt{/64});
\item введение нового downstream-клиента (например, новой
аналитической подсистемы);
\item появление дополнительных метаданных в маршрутных записях
(например, длины AS-path или признаков агрегации, см.
\cite{rfc4632}).
\end{itemize}

\subsection*{Г.5. Связь с теоретической моделью}

Эффект унификации источника, описанный выше, является частным
проявлением более общего архитектурного принципа: один производитель
истины, несколько потребителей, общий поток обновлений. Этот принцип
явно проводится в работах по log-oriented архитектурам
\cite{kreps2011kafka,kleppmann2017ddia}, в моделях stateful
streaming \cite{akidau2015dataflow,carbone2015flink} и в
наблюдаемости BGP-плана через BMP \cite{rfc7854}.

Таким образом, разработанный сервис не только решает прикладную
задачу, но и встраивается в более широкую инженерную традицию
проектирования систем работы с потоками изменений состояния.

\subsection*{Г.6. Возможные пути будущей интеграции}

Текущая интеграция сервиса с продуктами компании является базовой:
клиенты обращаются к нему по gRPC и используют возвращаемые ответы
напрямую. В будущем разумны следующие направления развития
интеграционного контура:

\begin{enumerate}
\item Кэширование на стороне клиента с инвалидацией через подписку
\texttt{SubscribePrefixASUpdates}.
\item Промежуточные адаптеры, преобразующие поток обновлений в
другие форматы (например, в очередь сообщений или в Kafka-топик
\cite{kreps2011kafka}).
\item Шардированное обслуживание, в котором несколько экземпляров
сервиса делят между собой нагрузку по AFI/SAFI или по диапазонам
ASN.
\item Многослойная топология, в которой продуктовые сервисы
подписываются на агрегированный поток событий, а не каждый раз
вычисляют lookup самостоятельно.
\item Интеграция с внешними route collectors \cite{ripe-ris,
routeviews} для верификации и обогащения локального состояния.
\end{enumerate}

\subsection*{Г.7. Итог приложения}

С точки зрения практической инженерии разработанный сервис
выполняет три функции одновременно:

\begin{itemize}
\item служит общей платформой для нескольких продуктов;
\item обеспечивает свежесть состояния на уровне реального времени;
\item даёт инфраструктурный фундамент для последующих расширений в
области сетевой аналитики и наблюдаемости.
\end{itemize}

В этом смысле работа имеет не только исследовательское, но и
непосредственное эксплуатационное значение, что подтверждает
заявленную во введении практическую значимость и сводит воедино
архитектурный и алгоритмический результаты диссертации.

